
\documentclass[12pt]{article}

% Layout.
\usepackage[top=1.2in, bottom=0.75in, left=1in, right=1in, headheight=1.0in, headsep=0pt]{geometry}

% Fonts.
\usepackage{mathptmx}
\usepackage[scaled=0.86]{helvet}
\renewcommand{\emph}[1]{\textsf{\textbf{#1}}}

% TiKZ.
\usepackage{tikz, pgfplots}
\usetikzlibrary{calc}
\pgfplotsset{my style/.append style={axis x line=middle, axis y line=middle, xlabel={$x$}, ylabel={$y$}}}
\pgfplotsset{compat=1.16}

% Misc packages.
\usepackage{amsmath,amssymb,latexsym,bm,array,multicol,enumitem}
\usepackage{graphicx}
\usepackage{xcolor}
\usepackage[most]{tcolorbox}
\usepgfplotslibrary{fillbetween}
\definecolor{solbg}{RGB}{225,240,252}
\definecolor{solframe}{RGB}{70,130,180}
\newtcolorbox{solution}[1][]{%
  breakable, enhanced,
  colback=solbg, colframe=solframe, boxrule=0.7pt, arc=3pt,
  left=6pt, right=6pt, top=5pt, bottom=5pt,
  before skip=6pt, after skip=6pt,
  fonttitle=\sffamily\bfseries\small, coltitle=white, title={Solution},
  #1}

% Commands to set various header/footer components.
\makeatletter
\def\doctitle#1{\gdef\@doctitle{#1}}
\doctitle{Use {\tt\textbackslash doctitle\{MY LABEL\}}.}
\def\docdate#1{\gdef\@docdate{#1}}
\docdate{Use {\tt\textbackslash docdate\{MY DATE\}}.}
\def\doccourse#1{\gdef\@doccourse{#1}}
\let\@doccourse\@empty
\def\docscoring#1{\gdef\@docscoring{#1}}
\let\@docscoring\@empty
\def\docversion#1{\gdef\@docversion{#1}}
\let\@docversion\@empty
\makeatother

% Headers and footers layout.
\makeatletter
\usepackage{fancyhdr}
\pagestyle{fancy}
\fancyhf{} % Clears all headers/footers.
\lhead{\emph{\@doctitle\hfill\@docdate} \medskip
\ifnum \value{page} > 1\relax\else\\
\emph{Name: \rule{3.5in}{1pt}\hfill \@docscoring}
\fi}

\rfoot{\emph{\@docversion}}
\lfoot{\emph{\@doccourse}}
\cfoot{\emph{\thepage}}
\renewcommand{\headrulewidth}{0pt}%
\makeatother

% Paragraph spacing
\parindent 0pt
\parskip 6pt plus 1pt

% A problem is a section-like command. Use \problem{5} for a problem worth 5 points.
\newcounter{probcount}
\newcounter{subprobcount}
\setcounter{probcount}{0}

\newcommand{\problem}[1]{%
\par
\addvspace{4pt}%
\setcounter{subprobcount}{0}%
\stepcounter{probcount}%
\makebox[0pt][r]{\emph{\arabic{probcount}.}\hskip1ex}\emph{[#1 points]}\hskip1ex}

\newcommand{\thesubproblem}{\emph{\alph{subprobcount}.}}

% Subproblems are an enumerate-like environment with a consistent
% numbering scheme.  Use \begin{subproblems}\item...\item...\end{subproblems}
\newenvironment{subproblems}{%
\begin{enumerate}[itemindent=8pt,leftmargin=0pt]%
\setcounter{enumi}{\value{subprobcount}}%
\renewcommand{\labelenumi}{\emph{\textsl{\alph{enumi})}} \,}%
}%
{%
\setcounter{subprobcount}{\value{enumi}}%
\end{enumerate}%
}

% Blanks for answers in normal and math mode.
\newcommand{\blank}[1]{\rule{#1}{0.75pt}}
\newcommand{\mblank}[1]{\underline{\hspace{#1}}}
\def\emptybox(#1,#2){\framebox{\parbox[c][#2]{#1}{\rule{0pt}{0pt}}}}

% Misc.
\renewcommand{\d}{\displaystyle}
\newcommand{\ds}{\displaystyle}
\newcommand{\threeopts}{{\small \hspace{-6mm} $\begin{matrix} \text{\textsc{converges}} \\ \text{\textsc{absolutely}} \end{matrix}$ \qquad\qquad $\begin{matrix} \text{\textsc{converges}} \\ \text{\textsc{conditionally}} \end{matrix}$ \qquad\qquad \textsc{diverges}} \bigskip}


\newcommand{\ba}{\mathbf{a}}
\newcommand{\bb}{\mathbf{b}}
\newcommand{\bc}{\mathbf{c}}
\newcommand{\bi}{\mathbf{i}}
\newcommand{\bj}{\mathbf{j}}
\newcommand{\bk}{\mathbf{k}}
\newcommand{\bu}{\mathbf{u}}
\newcommand{\bv}{\mathbf{v}}
\newcommand{\bw}{\mathbf{w}}

\newcommand{\ip}[2]{\mathrm{\left<#1,#2\right>}}
\newcommand{\vv}[2]{\mathrm{\left<#1,#2\right>}}
\newcommand{\vvv}[3]{\mathrm{\left<#1,#2,#3\right>}}


\doctitle{Math 252: Quiz 8 Solutions}
\docdate{5 November, 2026}
\doccourse{}
\docversion{Solutions}
\docscoring{\fbox{{\LARGE \strut}\hspace{0.8in} / 25}}

\begin{document}
30 minutes maximum.  No aids (book, calculator, etc.) are permitted.
Show all work and use proper notation for full credit.  Answers should
be in reasonably-simplified form.  25 points possible.

\begin{enumerate}

\item Determine whether the following series converge or diverge.
  \begin{enumerate}
  \item (5 points) $\ds \sum_{n=1}^{\infty} \left( \frac{1}{e} + \frac{1}{n} \right)^n$
    \begin{solution}
      The terms are positive, so we may use the root test.
      \[
        \lim_{n\to\infty} \sqrt[n]{\left( \frac1e + \frac1n \right)^n}
        = \lim_{n\to\infty} \left( \frac1e + \frac1n \right) = \frac1e < 1,
      \]
      so by the root test the series \boxed{\text{converges}}.
    \end{solution}
    \vfill
    
  \item (5 points) $\ds \sum_{n=100}^{\infty} \frac{(n!)^2}{(3n)!}$
    \begin{solution}
      The terms are positive, so we may use the ratio test.
      \begin{align*}
        \frac{a_{n+1}}{a_n}
          &= \frac{\big((n+1)!\big)^2}{(3n+3)!} \cdot \frac{(3n)!}{(n!)^2}
           = \frac{(n+1)^2}{(3n+3)(3n+2)(3n+1)}
           = \frac{n+1}{3(3n+2)(3n+1)} \\[4pt]
        \lim_{n\to\infty} \frac{a_{n+1}}{a_n} &= 0 < 1,
      \end{align*}
      so by the ratio test the series \boxed{\text{converges}}.
    \end{solution}
    \vfill
    
  \end{enumerate}

  \newpage


\item Consider the power series
  $\ds \sum_{n=0}^{\infty} \frac{(3x)^n}{2^n}$.
  \begin{enumerate}
  \item (5 points) Determine the radius of convergence of the power series.
    \begin{solution}
      $\ds\sum_{n=0}^{\infty} \frac{(3x)^n}{2^n} = \sum_{n=0}^{\infty} \left( \frac{3x}{2} \right)^n$
      is a geometric series with $r = \frac{3x}{2}$, so it converges exactly when
      \[
        \left| \frac{3x}{2} \right| < 1 \iff |x| < \frac23.
      \]
      \[
        \boxed{R = \frac23}
      \]
    \end{solution}
    \vfill
    
  \item (5 points) Determine the interval of convergence of the power series.
    \begin{solution}
      Check the endpoints.

      $x = \frac23$: $\ds\sum_{n=0}^\infty 1^n$; the terms do not approach $0$, so it diverges
      by the divergence test.

      $x = -\frac23$: $\ds\sum_{n=0}^\infty (-1)^n$; the terms do not approach $0$, so it
      diverges by the divergence test.
      \[
        \boxed{\left( -\frac23,\ \frac23 \right)}
      \]
    \end{solution}
    \vfill
  \end{enumerate}

  \newpage

\item (5 points) Use the fact that
  $\ds \frac{1}{1-x} = \sum_{n=0}^{\infty} x^n$ with interval of
  convergence $(-1,1)$ to find a power series representation for
  $\ds \frac{3}{1+5x}$.
  \begin{solution}
    Replace $x$ with $-5x$:
    \[
      \frac{3}{1+5x} = 3 \cdot \frac{1}{1-(-5x)} = 3\sum_{n=0}^{\infty} (-5x)^n
      = \boxed{\sum_{n=0}^{\infty} 3(-5)^n x^n}
    \]
    valid for $|-5x| < 1$, i.e.\ $-\frac15 < x < \frac15$.
  \end{solution}

  \vfill

\item (BONUS) (2 points) What is the center of the power series
  $\ds \sum_{n=0}^{\infty} \frac{3(x+5)^n}{4(n+1)}$? (This is not meant to
  be a challenging problem.)
  \begin{solution}
    $x + 5 = x - (-5)$, so the center is \boxed{-5}.
  \end{solution}

  \vfill
  

\end{enumerate}

\end{document}

%-------------------------------------------------------------------------------------------------------------------------------------------------------------------------------------------------------------------

%%% Local Variables:
%%% mode: latex
%%% TeX-master: t
%%% End:
